\documentclass[UTF8,11pt]{ctexart}
\usepackage[a4paper,margin=24mm]{geometry}
\usepackage{amsmath,amssymb,amsthm,mathtools,mathrsfs}
\usepackage[all]{xy}
\usepackage{xurl,hyperref,enumitem}
\usepackage{tikz}
\usetikzlibrary{arrows.meta,cd,decorations.markings,calc,patterns}
\hypersetup{hidelinks,breaklinks=true}
\setlength{\emergencystretch}{3em}
\allowdisplaybreaks
\newtheorem*{theorem}{Theorem}
\newtheorem*{thm}{Theorem}
\newtheorem*{lemma}{Lemma}
\newtheorem*{proposition}{Proposition}
\newtheorem*{prop}{Proposition}
\newtheorem*{corollary}{Corollary}
\newtheorem*{cor}{Corollary}
\newtheorem*{claim}{Claim}
\theoremstyle{definition}
\newtheorem*{definition}{Definition}
\newtheorem*{defn}{Definition}
\newtheorem*{remark}{Remark}
\newtheorem*{example}{Example}
\newcommand{\R}{\mathbb R}
\newcommand{\C}{\mathbb C}
\newcommand{\Z}{\mathbb Z}
\newcommand{\Q}{\mathbb Q}
\newcommand{\N}{\mathbb N}
\newcommand{\F}{\mathbb F}
\newcommand{\D}{\mathbb D}
\newcommand{\bB}{\mathbb B}
\newcommand{\bC}{\mathbb C}
\newcommand{\bR}{\mathbb R}
\newcommand{\bZ}{\mathbb Z}
\newcommand{\bN}{\mathbb N}
\newcommand{\bQ}{\mathbb Q}
\newcommand{\bD}{\mathbb D}
\newcommand{\bF}{\mathbb F}
\newcommand{\CP}{\mathbb{CP}}
\newcommand{\RP}{\mathbb{RP}}
\newcommand{\sO}{\mathscr O}
\newcommand{\sI}{\mathscr I}
\newcommand{\sF}{\mathscr F}
\newcommand{\sG}{\mathscr G}
\newcommand{\sH}{\mathscr H}
\newcommand{\sR}{\mathscr R}
\newcommand{\sM}{\mathscr M}
\newcommand{\sV}{\mathscr V}
\newcommand{\sU}{\mathscr U}
\DeclareMathOperator{\Hom}{Hom}
\DeclareMathOperator{\Ext}{Ext}
\DeclareMathOperator{\Tor}{Tor}
\DeclareMathOperator{\Spec}{Spec}
\DeclareMathOperator{\Proj}{Proj}
\DeclareMathOperator{\supp}{supp}
\DeclareMathOperator{\im}{im}
\DeclareMathOperator{\id}{id}
\DeclareMathOperator{\Tr}{Tr}
\DeclareMathOperator{\tr}{tr}
\DeclareMathOperator{\rank}{rank}
\DeclareMathOperator{\ord}{ord}
\DeclareMathOperator{\dist}{dist}
\DeclareMathOperator{\End}{End}
\DeclareMathOperator{\Aut}{Aut}
\DeclareMathOperator{\Gal}{Gal}
\DeclareMathOperator{\vol}{vol}
\DeclareMathOperator{\Cl}{Cl}
\DeclareMathOperator{\coker}{coker}
\DeclareMathOperator{\codim}{codim}
\DeclareMathOperator{\divergence}{div}
\DeclareMathOperator{\Ric}{Ric}
\DeclareMathOperator{\grad}{grad}
\DeclareMathOperator{\Hess}{Hess}
\newcommand{\abs}[1]{\left\lvert#1\right\rvert}
\newcommand{\sabs}[1]{\lvert#1\rvert}
\newcommand{\babs}[1]{\bigl\lvert#1\bigr\rvert}
\newcommand{\Babs}[1]{\Bigl\lvert#1\Bigr\rvert}
\newcommand{\bbabs}[1]{\biggl\lvert#1\biggr\rvert}
\newcommand{\BBabs}[1]{\Biggl\lvert#1\Biggr\rvert}
\newcommand{\norm}[1]{\left\lVert#1\right\rVert}
\newcommand{\snorm}[1]{\lVert#1\rVert}
\newcommand{\bnorm}[1]{\bigl\lVert#1\bigr\rVert}
\newcommand{\Bnorm}[1]{\Bigl\lVert#1\Bigr\rVert}
\newcommand{\bbnorm}[1]{\biggl\lVert#1\biggr\rVert}
\newcommand{\BBnorm}[1]{\Biggl\lVert#1\Biggr\rVert}
\newcommand{\widebar}[1]{\overline{#1}}
\newcommand{\nicefrac}[2]{\frac{#1}{#2}}
\newcommand{\linnprod}[2]{\langle#1,#2\rangle}
\newcommand{\myindex}[1]{#1}
\newcommand{\glsadd}[1]{}
\newcommand{\avoidbreak}{}
\newcommand{\sourceref}[1]{\textnormal{[原文：\nolinkurl{#1}]}}
\newcommand{\thmref}[1]{\sourceref{#1}}
\newcommand{\lemmaref}[1]{\sourceref{#1}}
\newcommand{\propref}[1]{\sourceref{#1}}
\newcommand{\corref}[1]{\sourceref{#1}}
\newcommand{\defnref}[1]{\sourceref{#1}}
\newcommand{\exampleref}[1]{\sourceref{#1}}
\newcommand{\exerciseref}[1]{\sourceref{#1}}
\newcommand{\figureref}[1]{\sourceref{#1}}
\newcommand{\sectionref}[1]{\sourceref{#1}}
\newcommand{\appendixref}[1]{\sourceref{#1}}
\newcommand{\stacksref}[2]{\href{https://stacks.math.columbia.edu/tag/#1}{#2 [#1]}}

\begin{document}
\section*{049\quad Young对称化子与对称群不可约表示的分类}
\subsection*{来源与计数目标}
Pavel Etingof 与 Oleg Golberg、Sebastian Hensel、Tiankai Liu、Alex Schwendner、Elena Yudovina、Dmitry Vaintrob 合著的 \emph{Introduction to Representation Theory} 课程讲义。MIT 18.712，Fall 2010 的 OCW 原讲义版本，不是限制转载的 2016 年 AMS 出版书。Chapter 4，\S4.12--4.13，Definition 4.35、Theorem 4.36、Lemmas 4.40--4.43 及完整证明，分章 PDF 第12--14页。。获取日期：2026-09-28。
\par\url{https://ocw.mit.edu/courses/18-712-introduction-to-representation-theory-fall-2010/84358595a02a73bced2c4e363a5d66f0_MIT18_712F10_ch4.pdf}
\par\url{https://ocw.mit.edu/courses/18-712-introduction-to-representation-theory-fall-2010/pages/lecture-notes/}
\par\textbf{本文件只计一个问题：}由 Young 表构造全部且互不同构的 $S_n$ 复不可约表示，证明构造、不可约性和穷尽性。
\subsection*{作者命题与人类证明}

\paragraph{4.12 Representations of $S_n$.}
In this subsection we give a description of the representations of the symmetric group $S_n$ for any $n$.
\begin{definition}[4.35]
A partition $\lambda$ of $n$ is a representation of $n$ in the form $n=\lambda_1+\lambda_2+\cdots+\lambda_p$, where $\lambda_i$ are positive integers, and $\lambda_i\geq\lambda_{i+1}$.
\end{definition}
To such $\lambda$ we will attach a Young diagram $Y_\lambda$, which is the union of rectangles $-i\leq y\leq-i+1$, $0\leq x\leq\lambda_i$ in the coordinate plane, for $i=1,\ldots,p$. Clearly, $Y_\lambda$ is a collection of $n$ unit squares. A Young tableau corresponding to $Y_\lambda$ is the result of filling the numbers $1,\ldots,n$ into the squares of $Y_\lambda$ in some way (without repetitions). For example, we will consider the Young tableau $T_\lambda$ obtained by filling in the numbers in the increasing order, left to right, top to bottom.

We can define two subgroups of $S_n$ corresponding to $T_\lambda$:
\begin{enumerate}
\item The row subgroup $P_\lambda$: the subgroup which maps every element of $\{1,\ldots,n\}$ into an element standing in the same row in $T_\lambda$.
\item The column subgroup $Q_\lambda$: the subgroup which maps every element of $\{1,\ldots,n\}$ into an element standing in the same column in $T_\lambda$.
\end{enumerate}
Clearly, $P_\lambda\cap Q_\lambda=\{1\}$. Define the Young projectors:
\[
a_\lambda:=\frac1{|P_\lambda|}\sum_{g\in P_\lambda}g,\qquad
b_\lambda:=\frac1{|Q_\lambda|}\sum_{g\in Q_\lambda}(-1)^g g,
\]
where $(-1)^g$ denotes the sign of the permutation $g$. Set $c_\lambda=a_\lambda b_\lambda$. Since $P_\lambda\cap Q_\lambda=\{1\}$, this element is nonzero.

The irreducible representations of $S_n$ are described by the following theorem.
\begin{theorem}[4.36]
The subspace $V_\lambda:=\C[S_n]c_\lambda$ of $\C[S_n]$ is an irreducible representation of $S_n$ under left multiplication. Every irreducible representation of $S_n$ is isomorphic to $V_\lambda$ for a unique $\lambda$.
\end{theorem}
The modules $V_\lambda$ are called the Specht modules.
\paragraph{4.13 Proof of Theorem 4.36.}
\begin{lemma}[4.40]
Let $x\in\C[S_n]$. Then $a_\lambda xb_\lambda=\ell_\lambda(x)c_\lambda$, where $\ell_\lambda$ is a linear function.
\end{lemma}
\begin{proof}
If $g\in P_\lambda Q_\lambda$, then $g$ has a unique representation as $pq$, $p\in P_\lambda$, $q\in Q_\lambda$, so $a_\lambda gb_\lambda=(-1)^q c_\lambda$. Thus, to prove the required statement, we need to show that if $g$ is a permutation which is not in $P_\lambda Q_\lambda$ then $a_\lambda gb_\lambda=0$.

To show this, it is sufficient to find a transposition $t$ such that $t\in P_\lambda$ and $g^{-1}tg\in Q_\lambda$; then
\[
a_\lambda gb_\lambda=a_\lambda tgb_\lambda
=a_\lambda g(g^{-1}tg)b_\lambda=-a_\lambda gb_\lambda,
\]
so $a_\lambda gb_\lambda=0$. In other words, we have to find two elements $i,j$ standing in the same row in the tableau $T=T_\lambda$, and in the same column in the tableau $T'=gT$ (where $gT$ is the tableau of the same shape as $T$ obtained by permuting the entries of $T$ by the permutation $g$). Thus, it suffices to show that if such a pair does not exist, then $g\in P_\lambda Q_\lambda$, i.e., there exists $p\in P_\lambda$, $q'\in Q'_\lambda:=gQ_\lambda g^{-1}$ such that $pT=q'T'$ (so that $g=pq^{-1}$, $q=g^{-1}q'g\in Q_\lambda$).

Any two elements in the first row of $T$ must be in different columns of $T'$, so there exists $q'_1\in Q'_\lambda$ which moves all these elements to the first row. So there is $p_1\in P_\lambda$ such that $p_1T$ and $q'_1T'$ have the same first row. Now do the same procedure with the second row, finding elements $p_2,q'_2$ such that $p_2p_1T$ and $q'_2q'_1T'$ have the same first two rows. Continuing so, we will construct the desired elements $p,q'$. The lemma is proved.
\end{proof}
Let us introduce the lexicographic ordering on partitions: $\lambda>\mu$ if the first nonvanishing $\lambda_i-\mu_i$ is positive.
\begin{lemma}[4.41]
If $\lambda>\mu$ then $a_\lambda\C[S_n]b_\mu=0$.
\end{lemma}
\begin{proof}
Similarly to the previous lemma, it suffices to show that for any $g\in S_n$ there exists a transposition $t\in P_\lambda$ such that $g^{-1}tg\in Q_\mu$. Let $T=T_\lambda$ and $T'=gT_\mu$. We claim that there are two integers which are in the same row of $T$ and the same column of $T'$. Indeed, if $\lambda_1>\mu_1$, this is clear by the pigeonhole principle (already for the first row). Otherwise, if $\lambda_1=\mu_1$, like in the proof of the previous lemma, we can find elements $p_1\in P_\lambda$, $q'_1\in gQ_\mu g^{-1}$ such that $p_1T$ and $q'_1T'$ have the same first row, and repeat the argument for the second row, and so on. Eventually, having done $i-1$ such steps, we'll have $\lambda_i>\mu_i$, which means that some two elements of the $i$-th row of the first tableau are in the same column of the second tableau, completing the proof.
\end{proof}
\begin{lemma}[4.42]
$c_\lambda$ is proportional to an idempotent. Namely, $c_\lambda^2=\frac{n!}{\dim V_\lambda}c_\lambda$.
\end{lemma}
\begin{proof}
Lemma 4.40 implies that $c_\lambda^2$ is proportional to $c_\lambda$. Also, it is easy to see that the trace of $c_\lambda$ in the regular representation is $n!$ (as the coefficient of the identity element in $c_\lambda$ is $1$). This implies the statement.
\end{proof}
\begin{lemma}[4.43]
Let $A$ be an algebra and $e$ be an idempotent in $A$. Then for any left $A$-module $M$, one has $\Hom_A(Ae,M)\cong eM$ (namely, $x\in eM$ corresponds to $f_x:Ae\to M$ given by $f_x(a)=ax$, $a\in Ae$).
\end{lemma}
\begin{proof}
Note that $1-e$ is also an idempotent in $A$. Thus the statement immediately follows from the fact that $\Hom_A(A,M)\cong M$ and the decomposition $A=Ae\oplus A(1-e)$.
\end{proof}
\begin{proof}[Completion of Theorem 4.36]
Now we are ready to prove Theorem 4.36. Let $\lambda\geq\mu$. Then by Lemmas 4.42, 4.43
\[
\Hom_{S_n}(V_\lambda,V_\mu)
=\Hom_{S_n}(\C[S_n]c_\lambda,\C[S_n]c_\mu)
=c_\lambda\C[S_n]c_\mu.
\]
The latter space is zero for $\lambda>\mu$ by Lemma 4.41, and $1$-dimensional if $\lambda=\mu$ by Lemmas 4.40 and 4.42. Therefore, $V_\lambda$ are irreducible, and $V_\lambda$ is not isomorphic to $V_\mu$ if $\lambda\ne\mu$. Since the number of partitions equals the number of conjugacy classes in $S_n$, the representations $V_\lambda$ exhaust all the irreducible representations of $S_n$. The theorem is proved.
\end{proof}

\subsection*{整理说明（非作者正文）}
保留原定义、两个组合消去论证、幂等元与 Hom 对应及分类末段；不把相关技术引理另计题目。标准先修是复有限群表示的完全可约性、Schur 引理和不可约表示数等于共轭类数；不能引用 Specht 分类本身。难点是从 Young 表的行列位置构造换位消去、由组合偏序控制群代数角空间，再返回表示分类。

原第12页把 $a_\lambda,b_\lambda$ 定义为归一化平均，但第14页 Lemma 4.42 的常数和单位元系数措辞使用未归一化写法；这两处均已查看原页，原文照录，不静默修订。所计目标是 4.36 的分类，不是 4.42 的精确常数；后续分类论证用到的是与非零幂等元成比例这一性质。此处记录可见记号不一致，不另行数学审稿。省略例子、练习及无关推论；三页原图均已取得。
\subsection*{可修改的简短标注（非作者正文）}
$c$：对称群不可约表示的显式分类。

$a$：Young 图与 Young 表；行列置换子群；对称化与反对称化；字典序；群代数幂等元；角空间与模同态；半单性。

\texttt{cross\_theory\_status = yes}。把 Young 表的组合行列约束转为群代数乘积的消去和一维角空间，再由幂等元的模同态对应得到不可约表示及其唯一参数化。位置：Lemmas 4.40--4.43 与 Theorem 4.36 末段。

全部标注均为非穷尽、可修改初稿；未标注不表示未使用。
\subsection*{许可与修改记录}
材料来自 MIT OpenCourseWare，作者与课程见来源。按该站所示 CC BY-NC-SA 4.0 许可复用，整理部分沿用同一许可。\url{https://creativecommons.org/licenses/by-nc-sa/4.0/}。2026-09-28：助手转录题证、调整数学排版并加入来源、范围和初稿标注；不暗示作者或 MIT 认可本次整理。所留为本次题证转录摘录，不是原 PDF 下载件。
\end{document}
